Whilst \textit{failures} denote the actual inability of a service to perform its functions \cite{Saltzer2009_PrinciplesComputerSystemDesign}, \textit{anomalies} correspond to the observable symptoms for such failures \cite{29_Wang2018_CloudRanger}, \eg a service slowing its response time, reducing its throughput, or logging error events.   
The problem of \textit{anomaly detection} in multi-service applications hence consists of identifying anomalies that can possibly correspond to failures affecting their services.
Anomalies can be detected either at \textit{application-level} or at \textit{service-level}, based on whether symptoms of a failure are observed by considering the application as whole (\eg performance degradations or errors shown by its frontend) or by focusing on specific services.  

\textit{Root cause analysis} then consists of identifying the reasons why an application- or service-level anomaly has been observed, with the ultimate goal of providing possible reasons for the corresponding failure to occur.
Root cause analysis is enacted with techniques analysing what has actually happened while a multi-service application was running.
\upd{
    It is hence not to be confused with debugging techniques, which look for the possible reasons for an observed anomaly by re-running applications in a testing environment and by trying to replicate the observed anomaly (\eg as in the delta debugging proposed by Zhou \etal \cite{03_Zhou2019_DeltaDebugging}).
    Root cause analysis techniques instead only analyse information collected while the application was running in production (\eg logs or monitoring data), without re-running the application itself. 
}

Anomaly detection and root cause analysis are typically enacted based on runtime information collected on the services forming an application, hereafter also called \textit{application services}. 
Such information includes \textit{KPIs} (Key Performance Indicators) monitored on application services, \eg response time, availability, or resource consumption, as well as the \textit{events} logged by application services.
Each logged event provide information on something happened to the service logging such event, \eg whether it was experiencing some error, hence logging an \textit{error event}, or whether it was invoking or being invoked by another service.
The events logged by a service constitute the \textit{service logs}, whereas the union of all \textit{service logs} distributed among the services forming an application is hereafter referred as \textit{application logs}.

Application logs are not to be confused with \textit{distributed traces}, obtained by instrumenting application services to feature \textit{distributed tracing}.
The latter is a distributed logging method used to profile and monitor multi-service applications, which essentially consists of (i) instrumenting the services forming an application to assign each user request a unique id, (ii) passing such id to all services that are involved in processing the request, (iii) including the request id in all log messages, and (iv) recording information (\eg start time, end time) about the service invocations and internal operations performed when handling a request in a service \cite{Richardson2018_MicroservicesPatterns}.
All the events generated by user requests are hence tracked in the distributed \enquote{traces}, and traces of events generated by the same request can be identified based on the request id. 